\chapter{Introduction}
\setcounter{equation}{0}
\section{Overview of URL Phishing}
Phishing is a social-engineering attack in which an attacker deceives a user into revealing sensitive information, transferring money, or approving an unintended action. A common first step is to distribute a deceptive Uniform Resource Locator (URL) that directs a victim to a fraudulent website or to a resource associated with malware. Such links may be shared through email, messaging applications, social networks, or compromised websites. Successful attacks can lead to financial loss, identity theft, and reputational damage. Malicious URLs may be associated with spam, phishing, malware distribution, or website defacement. Although these threats are not identical, many detection systems reduce the decision to a binary classification: benign or malicious.\cite{bib1}

\section{Motivation and Importance}
Browser-based protection needs to assess links quickly, including URLs that have not previously appeared in public blocklists. Static lists are valuable for known threats but cannot by themselves recognise every newly created or modified malicious URL. Machine-learning models can identify patterns in URL characteristics, but their performance may decline when the source or distribution of URLs changes. This is known as dataset shift or a generalisation gap.

The selected paper, by Misiek and Hyla, introduces X-PHIDE, an XGBoost-based phishing URL detection method and evaluates it across three datasets. Its central emphasis is not only high performance on a single test split, but also cross-dataset validation, feature stability, and practical computational cost.\cite{bib1}
